% !TeX root = ../教材题库_全部条目_转写工作稿.tex
% Source: XJ-B1-C04-SU-P150-Y0882
% 第4章 / 小结与复习 小结与复习 / 教材 p.150 / PDF 第157页
% OCR未校对：公式与文字均可能有误；未完成数学TeX转写；禁止直接发给学生
% 原题及相邻原解截图（核对用）：../题目原图/XJ-B1-C04-SU-P150-Y0882-01.webp
% 本文件只在转写完成后作为题干使用；OCR 条目尚不是可靠题干。
\begin{bankitem}{第4章 小结与复习 第2题}{第4章 / 小结与复习 小结与复习 / 教材 p.150 / PDF 第157页}{XJ-B1-C04-SU-P150-Y0882}{OCR未校对：公式与文字均可能有误；未完成数学TeX转写；禁止直接发给学生}
\begin{ocrdraft}
[教材p.150 / PDF第157页]
\par
2．有了幂运算，A一伊，在A，B，C三个数中，一个不变，另外两个数之
\par
间就有了确定的对应关系．其中一个作为自变量“另一个就是的函数．这
\par
就得到三类重要的函数：指数函数、对数函数和幂函数．如何理解对数运算是
\par
指数运算的逆运算？指数和对数有哪些运算法则？其中最基本的是什么？
\end{ocrdraft}
\end{bankitem}
